%%
%%% macros for Theories of Grammar
%%%
\usepackage{polyglossia}
\setdefaultlanguage{english}
%\setmainfont{TeX Gyre Pagella}


\newcommand{\logo}{~}
\newcommand{\header}[3]{%
\title{\vspace*{-2ex} \large JMORF/MORS --- Morpho-Syntax
\\[2ex] \Large  \emp{#2} \\ \emp{#3}}
\author{\blu{Francis Bond}   \\ 
\normalsize  \textbf{Palacký University}\\
\normalsize  \url{https://fcbond.github.io/}\\
\normalsize  \texttt{bond@ieee.org}}
\MyLogo{JMORF (2026)}
\renewcommand{\logo}{#2}
\hypersetup{
   pdfinfo={
     Author={Francis Bond},
     Title={#1: #2},
     Subject={JMORF/MORS: Morpho-Syntax},
     Keywords={Syntax, Morphology, HPSG, Unification, Constructions},
     License={CC BY 4.0}
   }
 }

\date{#1
  \\ Location: KB-1.11}
}
\usepackage{hyperref}
\hypersetup{
     colorlinks,
     linkcolor={blue!50!black},
     citecolor={red!50!black},
     urlcolor={blue!80!black}
   }



\usepackage{xcolor}
\usepackage{graphicx}
\newcommand{\blu}[1]{\textcolor{blue}{#1}}
\newcommand{\grn}[1]{\textcolor{green}{#1}}
\newcommand{\hide}[1]{\textcolor{white}{#1}}
\newcommand{\emp}[1]{\textcolor{red}{#1}}
\newcommand{\txx}[1]{\textbf{\textcolor{blue}{#1}}}
\newcommand{\lex}[1]{\textbf{\mtcitestyle{#1}}}


% \renewcommand{\labelitemi}{\textcolor{violet}{‣}}
% \renewcommand{\labelitemii}{\textcolor{purple}{‣}}

\newcommand{\subhead}[1]{\noindent\textbf{#1}\\[5mm]}

\newcommand{\Bad}{\emp{\raisebox{0.15ex}{\ensuremath{\mathbf{\otimes}}}}}
\newcommand{\bad}[1]{*\eng{#1}}

\newcommand{\com}[1]{\hfill (\emp{#1})}%

\usepackage{relsize,xspace}
\newcommand{\into}{\ensuremath{\rightarrow}\xspace}
\newcommand{\tot}{\ensuremath{\leftrightarrow}\xspace}
\usepackage{url}
\newcommand{\lurl}[1]{\MyLogo{\url{#1}}}

\usepackage{langsci-gb4e}
\let\eachwordone=\slshape
%%% langsci-gb4e defines \thexnumi as \@roman\@xsi{xnumi}, which only works
%%% inside an exe environment; some decks use xlisti on its own.  Restore
%%% the plain gb4e definition so that keeps working.
\makeatletter
\def\thexnumi{\@xsi{xnumi}}
\makeatother
\newcommand{\lx}[1]{\textbf{\mtciteform{#1}}}
\newcommand{\ix}{\ex\slshape}



\newcommand{\ent}{\ensuremath{\Rightarrow}\xspace}
\newcommand{\ngv}{\ensuremath{\not\Rightarrow}\xspace}
%\usepackage{times}
%\usepackage{nttfoilhead}
\newcommand{\myslide}[1]{\foilhead[-25mm]{\raisebox{12mm}[0mm]{\emp{#1}}}\MyLogo{\logo}}
\newcommand{\myslider}[1]{\rotatefoilhead[-25mm]{\raisebox{12mm}[0mm]{\emp{#1}}}}
%\newcommand{\myslider}[1]{\rotatefoilhead{\raisebox{-8mm}{\emp{#1}}}}

\newcommand{\section}[1]{\myslide{}{\begin{center}\Huge \emp{#1}\end{center}}}



\usepackage[lyons,j,e,k]{mtg2e}
%\renewcommand{\mtcitestyle}[1]{\textcolor{-red!75!green!50}{\textsl{#1}}}
\renewcommand{\mtcitestyle}[1]{\textcolor{teal}{\textsl{#1}}}
\newcommand{\cs}{\mtciteform}
\newcommand{\iz}[1]{\texttt{\textup{#1}}}
\newcommand{\gm}{\textsc}
\usepackage[normalem]{ulem}
\newcommand{\ul}{\uline}
\newcommand{\ull}{\uuline}
\newcommand{\wl}{\uwave}
\newcommand{\vs}{\ensuremath{\Leftrightarrow}~}
%%%
%%% Bibliography
%%%
\usepackage{natbib}
%\usepackage{url}
\usepackage{bibentry}



\usepackage[utf8]{inputenc}
%%% trees: forest is what we now recommend to students, but the older
%%% decks still draw their trees with rtrees/qtree.  Both are loaded
%%% while the decks are being migrated one at a time.
\usepackage{rtrees,qtree}
%%% rtrees defines \lf inside the tree environment, not at top level,
%%% so provide it first (cf. \val in headx.tex).
\providecommand{\lf}[1]{}
\renewcommand{\lf}[1]{\br{#1}{}}
\usepackage{multicol}
\usepackage{forest}
\useforestlibrary{linguistics}
\forestapplylibrarydefaults{linguistics}

%%% \makebox's optional arguments are read as AVM delimiters by
%%% langsci-avm, so hide them behind a macro with no visible brackets.
\newcommand{\lbox}[1]{\makebox[0mm][l]{#1}}

%%% Tree styles.
%%%   words  -- trees whose leaves are actual words (italicised, aligned)
%%%   cat    -- trees whose leaves are categories or feature structures
\forestset{
  words/.style={for tree={s sep=7mm, l sep=5mm,
                          if n children=0{font=\itshape, tier=terminal}{}}},
  cat/.style={for tree={s sep=7mm, l sep=5mm}},
}
\newcommand{\wtree}[1]{\begin{forest} words #1 \end{forest}}
\newcommand{\ctree}[1]{\begin{forest} cat #1 \end{forest}}
%%% the rule being applied at this step of a derivation
\newcommand{\step}[1]{\begin{center}\emp{#1}\end{center}}

%%% AVMs (Language Science Press)
%%% rtrees does \RequirePackage{avm}, and Manning's avm *environment*
%%% collides with langsci-avm's \avm *command*.  We only want rtrees for
%%% its tree macros, so retire the environment before loading langsci-avm.
%%% This whole block goes once the trees are migrated to forest.
\makeatletter
\let\avm\@undefined
\let\endavm\@undefined
\makeatother
\usepackage{langsci-avm}

%%% Macros the decks use, carried over from headx.tex
\newcommand{\blank}{\rule{3em}{1pt}\xspace}
\newcommand{\ft}[1]{\textsc{#1}}
\newcommand{\idn}{\mtciteform}    %%% Indonesian
\newcommand{\prd}[1]{\textbf{#1}}
\newcommand{\typ}[1]{\textit{#1}}
%%% \val may or may not be defined already, depending on the avm.sty
%%% that rtrees drags in; provide it first, then set our definition.
\providecommand{\val}[1]{}
\renewcommand{\val}[1]{\textit{#1}}
%%% pstricks oval nodes, used to ring parts of rtrees trees
\newcommand{\OV}[1]{\ovalnode[linestyle=dotted,linecolor=red]{A}{#1}}
\newcommand{\OVB}[1]{\ovalnode[linestyle=dotted,linecolor=blue]{A}{#1}}

%%% From CSLI book
\newcommand{\mc}{\multicolumn}
\newcommand{\HD}{\textbf{H}\xspace}
\newcommand{\el}{< >}       %%% empty list; langsci-avm needs the space

\newcommand{\A}{\noindent\textbf{A}: }
\newcommand{\Q}{\noindent\textbf{Q}: }
%\newcommand{\C}{\noindent\textbf{C}: }

